% Auto-generated by python/build_tables.py
\begin{table*}[!htbp]
\centering
\begin{threeparttable}
\caption{Interaction-selection robustness: agent-first versus edge-uniform activation.}
\label{tab:appendix-selection-rule-robustness}
\small
\begin{tabular}{p{0.18\textwidth}p{0.22\textwidth}rrrr}
\toprule
\textbf{Condition} & \textbf{Selection rule} & \textbf{Mean $T$} & \textbf{Median $T$} & \textbf{$P_{90}$} & \textbf{$P_{95}$} \\
\midrule
RB\_low & Agent-first & 7,277.6 & 6,981.8 & 9,837.3 & 10,825.6 \\
RB\_low & Edge-uniform & 6,264.6 & 5,972.0 & 8,481.5 & 9,413.8 \\
BS\_low & Agent-first & 8,436.8 & 7,991.6 & 11,414.2 & 12,794.1 \\
BS\_low & Edge-uniform & 6,309.8 & 5,983.2 & 8,596.5 & 9,448.5 \\
BS\_high & Agent-first & 4,531.4 & 4,457.9 & 5,470.8 & 5,781.3 \\
BS\_high & Edge-uniform & 3,383.2 & 3,348.1 & 4,029.7 & 4,281.1 \\
\bottomrule
\end{tabular}
\begin{tablenotes}
\footnotesize
\item Notes: Agent-first selection represents scarce actor-level interaction capacity. Edge-uniform selection activates ties uniformly and therefore weakens actor-level overload. Under edge-uniform selection, the concentration penalty in $BS_{low}$ largely disappears, while translation capability remains strongly beneficial.
\end{tablenotes}
\end{threeparttable}
\end{table*}
